\documentclass[10pt,a4paper]{amsart}
\usepackage[margin=19mm]{geometry}
\usepackage{amsmath,amssymb,mathtools}
\usepackage[scr=rsfs,cal=euler]{mathalfa}
\usepackage{xfrac}
\usepackage[unicode,hidelinks,bookmarks=false]{hyperref}
\newcommand{\ie}{i.e.\ }
\newcommand{\cf}{cf.\ }
\newcommand{\resp}{respectively\ }
\newcommand{\Subsection}{\subsection}
\providecommand{\nobreakdash}{\nobreak\hbox{-}}
\setlength{\emergencystretch}{2em}
\multlinegap0pt
\usepackage{amssymb}
\usepackage{mathtools}
\usepackage{xcolor}
\usepackage[normalem]{ulem}
\usepackage{bm}
\newtheorem{lemma}{Lemma}
\newtheorem{theorem}{Theorem}
\newtheorem{definition}{Definition}
\newtheorem{corollary}{Corollary}
\newtheorem{proposition}{Proposition}
\def\A{{\mathbb A}}
\def\F{{\mathbb F}}
\def\Ql{\overline{\mathbb{Q}}_\ell}
\DeclareMathOperator{\Tr}{Tr}
\def\Z{{\mathbb Z}}
\title{Factorization of Additive Polynomials and van der Geer--van der Vlugt curves in characteristic $2$}
\author{Tetsushi Ito, Daichi Takeuchi, Takahiro Tsushima}
\date{Source: arXiv:2605.16729v1 (CC BY 4.0)}
\begin{document}
\maketitle

\noindent\emph{Source and licence.} Factorization of Additive Polynomials and van der Geer--van der Vlugt curves in characteristic $2$. Version arXiv:2605.16729v1 (CC BY 4.0), distributed by arXiv under the Creative Commons Attribution 4.0 International licence, http://creativecommons.org/licenses/by/4.0/, as recorded in the arXiv licence metadata for this exact version and shown on the abstract page https://arxiv.org/abs/2605.16729v1. Full licence notice: \texttt{licenses/arxiv-2605.16729-CC-BY-4.0.txt}.\par\smallskip

\noindent\textit{Editorial scope.} Counted unit: the article's proposition period (source0.tex lines 2462--2478). The article's complete original proof of it is delivered with it (source0.tex lines 2479--2558, one \texttt{proof} environment of 80 lines). No mathematics is written, repaired or re-derived: the statement and the proof are the article's own lines, in the article's own notation, and no retained line is substituted or re-flowed.  Retained uncounted as the source-local context the proof needs: lines 491--510; lines 1074--1076; lines 1251--1253; lines 1338--1361; lines 1706--1711; lines 1720--1727; lines 1758--1766; lines 2073--2082; lines 2277--2279; lines 2290--2305.  Nothing was omitted from any retained span, so the package carries no omission record.  No retained line contains a citation, so the wrapper carries no bibliography and no bibliography entry.  No retained line contains a figure environment, an image-inclusion command or drawing code, so no image is delivered and the package declares no asset dependency.  Editorial formatting only, in addition: an AMS \texttt{amsart} preamble that restates the article's own declarations and macros used by the retained lines, namely the environments \texttt{lemma}, \texttt{theorem}, \texttt{definition}, \texttt{corollary}, \texttt{proposition} on separate counters, together with the standard packages those lines need.  The retained environments print in this document as Lemma 1, Theorem 1, Definition 1, Corollary 1, Proposition 1 in order of appearance, whereas the article prints the same items with the numbering of its own full document.  The complete native source of the article as distributed in the arXiv e-print for this exact version is bundled unchanged as \texttt{sources/200743/source0.tex}.
\begin{lemma}\label{basic}
The following statements hold: 
\begin{enumerate}
\item The map 
\[\overline{k}^\times\times\mathbb{Z}\times{\mathcal V}\to \overline{k}\{\tau^{\pm1}\}\setminus\{0\},\qquad (a,n,W) \to a\tau^nf_W \]
 is bijective. In particular, for every nonzero $f\in \overline{k}\{\tau^{\pm1}\}$, there exists a unique triple $(a,n,W)$ such that 
 \[
 f=a\tau^nf_W. 
 \]
 Moreover, we have $\ker f=W$. 
 
\item Let $f=\sum_{i=r}^sa_i\tau^i$ with $a_ra_s\neq0$. Then 
\[\dim_{\F_p}\ker f=s-r. \]
\item For $f,g\in\overline{k}\{\tau^{\pm1}\}$, the following conditions are equivalent: 
\begin{enumerate}
\item There exists $h\in\overline{k}\{\tau^{\pm1}\}$ such that $f=hg$. 
\item We have $\ker g\subset \ker f$. 
\end{enumerate}
\end{enumerate}
\end{lemma}

 \begin{equation}\label{Fdef}
 F=\sum_{i=0}^eb_i\tau^i \in \F_q\{\tau^{\pm1}\}
 \end{equation}

\begin{equation}\label{Cadf}
D_{F,t}\colon y^p - y = x R_{F,t}(x), 
\end{equation}

 \begin{theorem}\label{FEV1}
 Assume that $F$ satisfies conditions \textup{(i)--(iii)} stated after \eqref{Fdef}. 
 
 Let $t\in\F_q$, and consider the curve $D_{F,t}$ defined by \eqref{Cadf}. 
  \begin{enumerate}
 \item 
 There exists an ${\rm Fr}_q$-equivariant isomorphism 
 \[H^1_c(D_{F,t,\F},\Ql)\cong \bigoplus_{v\in W_F^\ast}H^1_c(\A^1_{\F},{\mathcal L}_{\xi_p}(x,(t+v)x))^{\oplus p-1}. \]
 \item Each cohomology group $H^1_c(\A^1_{\F},{\mathcal L}_{\xi_p}(x,(t+v)x))$ has dimension $1$. 
 \item Let $\tau_v$ denote the ${\rm Fr}_q$-eigenvalue on  $H^1_c(\A^1_{\F},{\mathcal L}_{\xi_p}(x,(t+v)x))$.  Then 
 \[\tau_v=Q_q(t+v)^{-1}(-1-\sqrt{-1})^{[\F_q:\F_2]}.\]
 \item 
  Let $\overline{D}_{F,t}$ denote the smooth compactification of $D_{F,t}$. Then the $L$-polynomial 
\[
L(\overline{D}_{F,t}/\F_q,T):=\det(1-{\rm Fr}_qT \mid H^1(\overline{D}_{F,t,\F},\Ql))
\]
satisfies 
\[
L(\overline{D}_{F,t}/\F_q,T)=
\prod_{v\in W_F^\ast}(1-\tau_v T)^{p-1},
\]
where $\tau_v=Q_q(t+v)^{-1}(-1-\sqrt{-1})^{[\F_q:\F_2]}$. 
 \end{enumerate}
 \end{theorem}

\begin{definition}\label{mmx}
Let $\overline{C}$ be a geometrically connected smooth projective curve of genus $g$ over a finite field $k$. Let $L/k$ be  a finite field extension with $|L|=r$. Then  the Hasse--Weil bound gives 
\[1+r-2g\sqrt{r}\leq|\overline{C}(L)|\leq 1+r+2g\sqrt{r}.\]


We say that $\overline{C}$ is  \emph{$L$-maximal} (resp.\ \emph{$L$-minimal}) if the upper bound (resp.\ the lower bound) is achieved. We say that $\overline{C}$ is \emph{$L$-extremal} if it is either $L$-maximal or $L$-minimal. \end{definition}

\begin{corollary}\label{exN}
 If there exists $a\in \F_q$ such that $C_a$ is $\F_q$-extremal, then 
 \[
 V_q=V. 
 \]
  In other words, we have $V\subset \F_q$. 
 %Assume that $V\subset \F_q$. Then $C_a$ is $\F_q$-maximal (resp.\ $\F_q$-minimal) if and only if $a\in T_{R,\max}$ (resp. $a\in T_{R,\min})$. 
\end{corollary}

 \begin{lemma}\label{even}
 Let 
 \[R(x)=\sum_{i=1}^e a_i x^{p^i} \in 
 \F_q[x],\] 
 where $e\geq1$ and $a_e\neq0$. Set 
 \[
 E:=R+R^\ast,\qquad V:=\ker(E\colon \F\to \F).\]
 If $V\subset \F_q$, then $[\F_q:\F_p]$ is even. 
 \end{lemma}

\begin{lemma}\label{Trt_0}
Let $r=2^s$ be a power of $2$. 
Let $t_0\in\F$ be an element satisfying $t_0^{r}+t_0=1$. The following hold: 
\begin{enumerate}
\item We have $t_0\in\F_{r^2}$ and $t_0^{r+1}\in \F_{r}$. 
\item We have $\Tr_{r/2}(t_0^{r+1})=1$. 
\item Let $\Tr_{r^2/2}\colon W_2(\F_{r^2})\to W_2(\F_2)$ be the trace map. Then 
\[\Tr_{r^2/2}(t_0,0)=s(1,0)+(0,1). \]
\end{enumerate}
\end{lemma}

\begin{equation}\label{xpax}
X_a\colon y^p+y=x^{p+1}+ax^2
\end{equation} 

\begin{proposition}\label{Hermitian curve}
Let $p$ be a power of $2$, and let $q=p^m$ be a power of $p$.   We assume that $m$ is even, and write $m=2m_0$. 

Let $a\in \F_q$, and define $\alpha:=\Tr_{q/p^2}(a)$.  Consider the curve $X_a$ defined in \eqref{xpax}. Then the following hold: 
\begin{enumerate}
\item If $\alpha\neq0$, then $X_a$ is not $\F_q$-extremal (in the sense of Definition \ref{mmx}). The ${\rm Fr}_q$-eigenvalues are 
\[\pm(\sqrt{-1})^{t(\alpha^{p+1})}2^{[\F_q:\F_2]/2}, \]
each with the same multiplicity. Here, $t(\alpha^{p+1})$ denotes the integer in $\{0,1\}$ congruent to $\Tr_{p/2}(\alpha^{p+1})$ modulo $2$. 
\item If $\alpha=0$, then the curve $X_a$ is $\F_q$-extremal. Furthermore, $X_a$ is $\F_q$-maximal if and only if 
\[m_0+\Tr_{p/2}(\beta)=1,\]
 where $\beta\in\F_p$ is given by 
\[
\beta:=\sum_{\substack{0\leq i<j\leq m-1 \\ 2\mid j,\ 2\nmid i}}a^{p^i}a^{p^j}. 
\]
\end{enumerate}
\end{proposition}

\begin{proposition} \label{not well known case for period}
Assume that $\mu=2m$ where $m$ is even. Then, in each of the following cases, the given $R(x)$ satisfies 
\[
(\mu,1)=(\mu(\overline{C}_R),\delta(\overline{C}_R)). 
\]
\begin{enumerate}
    \item If $m=2$ and $\F_p\neq\F_2$, then 
    \[
    R(x):=x^p+ax, 
    \]
    where $a\in \F_p^\times$ satisfies $\Tr_{p/2}(a)=0$. 
\item If $m\geq4$, then 
    \[
R(x):=x^p+x^{p^3}+\cdots+x^{p^{m-1}}=\sum_{i=1}^{m/2}x^{p^{2i-1}}. 
\]
\end{enumerate}
\end{proposition}

\begin{proof}
(i) In this case, by Proposition~\ref{Hermitian curve}(i) applied to $q=p^2$, the curve $C_R=X_a$ is not $\F_{p^2}$-extremal, and the ${\rm Fr}_{p^2}$-eigenvalues are given by 
\[
\pm 2^{[\F_{p}:\F_2]}. 
\]
Therefore, $C_R$ is $\F_{p^4}$-minimal. This proves 
\[
(\mu(\overline{C}_R),\delta(\overline{C}_R))=(4,1). 
\]

\medskip

\noindent
(ii) In this case, define 
\[
E:=R+R^\ast,\qquad V:=\ker E. 
\]
Let $N$ be the smallest integer $\geq1$ such that 
\[
V\subset \F_{p^N}. 
\]
We first claim that $\mu\leq N$. Indeed, by comparing dimensions, we have 
\[
2(m-1)=\dim_{\F_p}
V\leq N.\]
On the other hand, by Lemma~\ref{even}, $N$ must be even. Hence, it suffices to show that $2(m-1)<N$. 

Suppose, to the contrary, that $2(m-1)=N$. Then, we must have $V=\F_{p^N}$. By Lemma~\ref{basic}(iii), it follows that 
\[
E=a\tau^M(\tau^N+1)
\]
for some $a\in \F^\times$ and $M\in\Z$. However, this is impossible, since $E$ has more than two nonzero coefficients. 
Thus, we conclude that 
\[\mu=2m\leq N. \]

By Corollary~\ref{exN}, for every integer $n\geq1$, extremality of $C_R$ over $\F_{p^n}$ implies that $V\subset \F_{p^n}$. Hence, to prove the assertion, it suffices to show that $C_R$ is $\F_{p^\mu}$-minimal. 

\medskip 

We show that $C_R$ is $\F_{p^\mu}$-minimal. Set $e:=m-1$ and $q:=p^\mu$. Define 
    \[
    F:=\sum_{i=0}^{m-1}\tau^i\in\F_p\{\tau^{\pm1}\}. 
    \]
    Since $m$ is even, we have $F^\ast(1)=0$. We claim that 
    \begin{equation}\label{WF=kerTr}
        W_F^\ast:=\ker F^\ast=\ker(\Tr_{p^m/p}\colon \F_{p^m}\to \F_p).
    \end{equation}
Indeed, the inclusion $\supset$ is clear from a direct computation, and then the equality follows by comparing cardinalities. 

Consequently, the element $F$ satisfies conditions~\textup{(i)--(iii)} stated after \eqref{Fdef}. Observe that 
    \[
    R+R^\ast=F^\ast F. 
    \]
Moreover, 
\[
\alpha_F(t)=\frac{m(m-1)}{2}+F^\ast(t)^2. 
\]
    Let $t_0\in\F_{p^2}$ be such that $t_0^{p}+t_0=1$. Then 
    \[
    F^\ast(t_0)=\sum_{i=0}^{m-1}t_0^{p^{-i}}=\sum_{i=0}^{m-1}(t_0-i)=\frac{m(m-1)}{2}, 
    \]
    and hence $\alpha_F(t_0)=0$. Thus, we obtain $D_{F,t_0}=C_R$. 

Therefore, by Theorem~\ref{FEV1}(iv), the ${\rm Fr}_q$-eigenvalues of $\overline{C}_R$ are given by 
\[
Q_q(t_0+v)^{-1}(\sqrt{-1})^{[\F_q:\F_2]/2}\cdot 2^{[\F_q:\F_2]/2}\qquad (v\in W_F^\ast). 
\]
It remains to show that $Q_q(t_0+v)=(\sqrt{-1})^{[\F_q:\F_2]/2}$ for all $v\in W_F^\ast$. 
We compute 
\[
Q_q(t_0+v)=\psi_q(t_0v)Q_q(t_0)Q_q(v). 
\]
By \eqref{WF=kerTr} and $t_0\in\F_{p^2}$, it follows that  
\begin{align*}
\psi_q(t_0v)&=1,\\
    Q_q(t_0)&=Q_{p^2}(t_0)^m=\bigl(-(\sqrt{-1})^{[\F_p:\F_2]}\bigr)^m,\qquad(\text{Lemma~\ref{Trt_0}(iii)})\\
    Q_q(v)&=\xi_{p^m}(\Tr_{p^{2m}/p^m}(v,0))=\xi_{p^m}(0,v^2)=\psi_p\bigl(\Tr_{p^m/p}(v)^2\bigr)=1. 
\end{align*}
Since $m$ is even, we obtain $Q_q(t_0+v)=(\sqrt{-1})^{[\F_q:\F_2]/2}$, and the curve $C_R$ is $\F_{p^\mu}$-minimal. This completes the proof. 
\end{proof}

\end{document}
